\documentclass[11pt,a4paper]{article}
\usepackage[T1]{fontenc}
\usepackage[utf8]{inputenc}
\usepackage{lmodern}
\usepackage[margin=2.5cm]{geometry}
\usepackage{microtype}
\usepackage{amsmath}
\usepackage{booktabs}
\usepackage{enumitem}
\usepackage[hidelinks]{hyperref}
\setlist{itemsep=1pt,topsep=3pt,parsep=0pt}

\title{\textbf{Sonifying the Heart: Inner Emotion Detection\\ from ECG through Music}\\[4pt]
  \large Music Workshop Sprint 3}
\author{%
  \begin{tabular}{c@{\hspace{1.5cm}}c@{\hspace{1.5cm}}c}
    Chervith Reddy & Manidhar Sukasi & Sathwik Reddy\\
    \texttt{2024101076} & \texttt{2023101067} & \texttt{2024121002}
  \end{tabular}}
\date{}

\begin{document}
\maketitle
\vspace{-1.5em}

\section{The Project}
In Sprint~1 we proposed turning ECG recordings into music, so that a healthy and an irregular
heart rhythm sound different. In Sprint~2 we planned the pipeline.

In Sprint~3 we built it, and added a fourth stage. That stage reads the emotion of the generated
music, as a window into the person's inner state:
\begin{center}
  \textbf{ECG data} $\;\rightarrow\;$ \textbf{Heartbeats} $\;\rightarrow\;$ \textbf{Music}
  $\;\rightarrow\;$ \textbf{Emotion}
\end{center}
The code is a Python package, \texttt{ecgmusic}, with a step-by-step notebook,
\texttt{main.ipynb}. One command turns a PhysioNet record, or any ECG file, into MIDI, audio, a
score plot and an emotion estimate.

\section{What We Built}
\begin{itemize}
  \item \textbf{ECG data.}
    \begin{itemize}
      \item The main source is the MIT-BIH Atrial Fibrillation Database (\texttt{afdb}). Most of
        its recordings contain both normal rhythm and AFib, so the same person is heard in both
        states.
      \item \texttt{mitdb} adds records 201 and 202, which come from one patient.
      \item Records are streamed from PhysioNet. Any WFDB or CSV file also works.
    \end{itemize}
  \item \textbf{Heartbeats.}
    \begin{itemize}
      \item R-peaks come from the annotations, or from the XQRS detector for other files. From
        them we compute the R--R intervals and RMSSD.
      \item Each recording gives a 30\,s normal clip (the steadiest window) and a 30\,s AFib clip
        (the best-detected window). We choose windows this way because \texttt{afdb}'s automatic
        beat marks jitter in places.
    \end{itemize}
  \item \textbf{Music.}
    \begin{itemize}
      \item \emph{Notes:} one per heartbeat. Each starts at the R-peak, lasts 85\% of the R--R
        interval, and is as loud as the ECG amplitude at that beat.
      \item \emph{Pitch:} follows the deviation $d_i$ of the interval from the mean of the last
        five intervals. It is C4 when the beat is on time, higher for a faster beat and lower
        for a slower one.
      \item \emph{Scale:} C-major pentatonic while $|d_i|<12\%$, chromatic beyond that.
      \item \emph{Harmony:} strings play C--Am--F--G. Each chord becomes diminished
        ($\geq 12\%$) or a cluster ($\geq 25\%$) as the beats get irregular.
      \item \emph{Pulse:} a soft ``lub-dub'' drum beat follows every heartbeat.
    \end{itemize}
  \item \textbf{Emotion.}
    \begin{itemize}
      \item Our model places the melody on Russell's valence--arousal plane, using three cues:
        rhythmic irregularity, pitch spread and consonance.
      \item The emotion picks the melody instrument: nylon guitar for calm, distortion guitar for
        tense, marimba for happy and cello for sad.
      \item Two pretrained models, Essentia and music2emo, give independent opinions.
    \end{itemize}
\end{itemize}

\section{Results}
We made 14 pieces: a normal and an AFib clip from each of 7 recordings of 6 people
(Table~\ref{tab:results}).\footnote{These numbers come from our first implementation. The
\texttt{ecgmusic} package uses the same rules and parameters, and will regenerate them.}

\begin{table}[h]
\centering
\small
\setlength{\tabcolsep}{4pt}
\caption{Every clip. Chords: consonant / tense / chaotic. Emotion from our model.}
\label{tab:results}
\begin{tabular}{llrrcl}
\toprule
Recording & Rhythm & RMSSD (ms) & Off-scale notes & Chords & Emotion (valence, arousal)\\
\midrule
afdb 04015 & normal &   7.6 &  0\% & 10 / 0 / 0 & calm ($+1.00$, $-0.97$)\\
           & AFib   & 126.1 & 57\% &  5 / 9 / 3 & tense ($-0.67$, $+0.76$)\\
afdb 04043 & normal &   8.5 &  0\% & 14 / 0 / 0 & calm ($+1.00$, $-0.96$)\\
           & AFib   & 134.9 & 52\% & 6 / 10 / 0 & tense ($-0.51$, $+0.45$)\\
afdb 04126 & normal &  18.6 &  0\% & 10 / 0 / 0 & calm ($+0.93$, $-0.79$)\\
           & AFib   & 141.9 & 66\% & 3 / 12 / 5 & tense ($-0.89$, $+0.85$)\\
afdb 06995 & normal &  17.8 &  0\% & 10 / 0 / 0 & calm ($+0.93$, $-0.81$)\\
           & AFib   & 122.8 & 51\% &  6 / 7 / 1 & tense ($-0.45$, $+0.31$)\\
afdb 08455 & normal &   8.3 &  0\% & 10 / 0 / 0 & calm ($+1.00$, $-0.96$)\\
           & AFib   &  99.9 & 31\% & 11 / 4 / 0 & happy, borderline ($+0.07$, $+0.02$)\\
mitdb 201  & normal &  64.1 &  0\% &  6 / 0 / 0 & calm ($+0.89$, $-0.69$)\\
           & AFib   & 194.3 & 52\% &  3 / 9 / 1 & tense ($-0.58$, $+0.76$)\\
mitdb 202  & normal &  37.4 &  0\% &  7 / 0 / 0 & calm ($+0.89$, $-0.65$)\\
           & AFib   & 146.9 & 33\% &  9 / 5 / 2 & tense ($-0.07$, $+0.44$)\\
\bottomrule
\end{tabular}
\end{table}

\begin{itemize}
  \item \textbf{Heartbeats.} In every recording, the AFib clip is 3--17 times as variable as the
    normal clip.
  \item \textbf{Music.}
    \begin{itemize}
      \item Normal pieces stay on the scale, and all 67 of their chords are consonant.
      \item AFib pieces leave the scale for 31--66\% of their notes, and 68 of their 111 chords
        turn tense or chaotic.
    \end{itemize}
  \item \textbf{Emotion.} Our model reads all 7 normal pieces as calm, and 6 of the 7 AFib pieces
    as tense. The seventh, 08455, is a mild AFib stretch and lands at the centre of the plane.
  \item \textbf{Pretrained models.} Both hear a difference in every recording, but describe it
    differently:
    \begin{itemize}
      \item Essentia hears the AFib pieces as more energetic and more pleasant (7 of 7).
      \item music2emo hears them as less pleasant and less energetic (6 of 7).
    \end{itemize}
    Tempo probably explains part of this. Every AFib clip is also the faster heart (101--160
    against 50--108\,bpm), and our model ignores tempo on purpose.
\end{itemize}

\section{Changes from the Sprint 2 Plan}
\begin{itemize}
  \item \textbf{Data source.} \texttt{afdb} is the main source, instead of \texttt{mitdb} with
    \texttt{nsrdb} as the reference. This lets us hear the same person in both rhythms.
  \item \textbf{Clip selection.} Choosing the steadiest window replaced the 0.3--2.0\,s
    plausibility filter. It handles the jittery beat marks we found.
  \item \textbf{Scale and loudness.}
    \begin{itemize}
      \item The scale is C-major pentatonic, which has the same notes as A minor pentatonic, to
        fit the C--Am--F--G harmony.
      \item The ECG amplitude sets the loudness. $|d_i|$ drives the harmony instead.
    \end{itemize}
  \item \textbf{Added:} the three-part arrangement, the emotion stage, and support for any ECG
    file.
  \item \textbf{Not done yet:} SDNN, pNN50, Poincar\'e plots and rate-matched clips.
\end{itemize}

\section{Next Steps}
\begin{itemize}
  \item Regenerate the results with the \texttt{ecgmusic} package.
  \item Test whether the emotion matches what people actually feel. We plan to use ECG datasets
    with self-reported emotion (WESAD, DREAMER, AMIGOS) and a listening study.
  \item Adapt the mapping for healthy hearts. There, calm breathing \emph{raises} beat-to-beat
    variability, which our current mapping would read as tension.
\end{itemize}

\section{Contributions}
\begin{tabular}{@{}lll@{}}
Chervith Reddy  & \texttt{2024101076} & Stages 1--2: ECG data and heartbeats\\
Manidhar Sukasi & \texttt{2023101067} & Stage 3: music, and integration (package, notebook)\\
Sathwik Reddy   & \texttt{2024121002} & Stage 4: emotion and evaluation\\
\end{tabular}

\end{document}
